\section*{Chapitre \ref{ch:b3:quantum-angular-momentum} --- Le moment cinétique quantique}

\begin{solution}{exo:b3:quantum-angular-momentum:1}
(a) $[\hat y\hat p_z - \hat z\hat p_y,\ \hat z\hat p_x - \hat x\hat
p_z]$~: seuls survivent les termes partageant une paire $\hat z, \hat p_z$,
ce qui donne $\iu\hbar(\hat x\hat p_y - \hat y\hat p_x)$. (b) $[\hat
L^2, \hat L_z] = \sum_i[\hat L_i^2, \hat L_z] = \hat L_x[\hat L_x,
\hat L_z] + [\hat L_x, \hat L_z]\hat L_x + (x \to y)$~: les quatre
termes se compensent deux à deux. (c) $[\hat L_z, \hat x] = \iu\hbar\hat y$~:
$\hat L_z$ engendre les rotations autour de $z$ --- des opérateurs
comme des états. (d) Que deux composantes soient simultanément bien
définies force (par le \omterm{def:b3:quantum-formalism:commutator}{commutateur}) la troisième à l'être aussi, avec
la valeur zéro pour les trois~: seul $l = 0$ y parvient.
\end{solution}

\begin{solution}{exo:b3:quantum-angular-momentum:2}
(a) $\hat L_z = \hbar\operatorname{diag}(1, 0, -1)$. (b) $\hat L_+ =
\hbar\sqrt2\,\begin{pmatrix}0&1&0\\0&0&1\\0&0&0
\end{pmatrix}$, et $\hat L_-$ sa transposée. (c) $\hat L_x =
\dfrac{\hbar}{\sqrt2}\begin{pmatrix}0&1&0\\1&0&1\\0&1&0
\end{pmatrix}$~: polynôme caractéristique $\lambda(\lambda^2 -
\hbar^2)$. (d) Le vecteur propre $L_x = +\hbar$ est $\tfrac12(1,
\sqrt2, 1)$~: probabilités $\tfrac14, \tfrac12, \tfrac14$ pour $m =
+1, 0, -1$.
\end{solution}

\begin{solution}{exo:b3:quantum-angular-momentum:3}
(a) $I = \mu r_0^2 = \qty{1.46e-46}{kg.m^2}$~; $B = \hbar^2/2I =
\qty{3.8e-23}{J} = \qty{0.24}{meV}$. (b) $2B/h = \qty{115}{GHz}$,
$\lambda = \qty{2.6}{mm}$. (c) \qty{231}{GHz} et \qty{346}{GHz}.
(d) L'espacement donne $B$, donc $I = \mu r_0^2$~: la longueur de
liaison d'une molécule située à des années-lumière, à peu près à la
moitié de la précision relative de l'espacement.
\end{solution}

\begin{solution}{exo:b3:quantum-angular-momentum:4}
(a) $m\hbar$ avec $m = -2 \dots 2$. (b) $\cos\theta = 2/\sqrt6$~:
$\theta = 35.3^\circ$. (c) Un alignement parfait rendrait $L_x = L_y
= 0$ bien définis en même temps que $L_z$ --- ce qu'interdisent les
\omterm{def:b3:quantum-formalism:commutator}{commutateurs}, sauf dans le cas trivial $l = 0$. (d) $l/\sqrt{l(l+1)} \to 1$~: aux
grands $l$ le cône se referme sur l'axe et la flèche classique
revient.
\end{solution}

\begin{solution}{exo:b3:quantum-angular-momentum:5}
(a) Sans dépendance en $\varphi$, l'opérateur donne
$-\hbar^2(\sin\theta)^{-1}\partial_\theta(\sin\theta\,(-\sin\theta))
= 2\hbar^2\cos\theta$. (b) Le même calcul avec le facteur
$\eu^{\pm\iu\varphi}$~: \omterm{def:b3:quantum-formalism:observable}{valeur propre} $2\hbar^2$~; $\hat L_z$
donne $\pm\hbar$. (c) $\cos\theta$ et $\sin\theta\,\eu^{\pm\iu
\varphi}$ changent de signe sous $\vect r \to -\vect r$ ($\theta \to \pi
- \theta$, $\varphi \to \varphi + \pi$)~: parité $-1 = (-1)^1$. (d)
Ce sont des vecteurs propres de l'\omterm{def:b3:quantum-formalism:observable}{opérateur hermitien} $\hat L_z$ pour
des \omterm{def:b3:quantum-formalism:observable}{valeurs propres} différentes.
\end{solution}

\begin{solution}{exo:b3:quantum-angular-momentum:6}
(a) $2J + 1$ états partagent le niveau (dégénérescence)~; le facteur
de Boltzmann taxe son énergie. (b) Maximiser le produit~:
$\dd/\dd J = 0$ donne $2 = (2J+1)^2B/k_{\text{B}}T$, soit le
$J_{\max}$ annoncé. (c) À \qty{10}{K}~: $J_{\max} \approx 0.8$, donc
$J = 1$ domine et les raies à \qty{115}{GHz} et \qty{230}{GHz}
brillent~; à \qty{300}{K}~: $J_{\max} \approx 7$. (d) $T \approx 2B(J_{\max} + \tfrac12)^2/
k_{\text{B}} \approx \qty{310}{K}$ --- un thermomètre rotationnel.
\end{solution}

\begin{solution}{exo:b3:quantum-angular-momentum:7}
(a) Le niveau supérieur se scinde en cinq, l'inférieur en trois, tous
au même espacement $\mu_{\text{B}}B$. (b) Les espacements étant égaux,
$h\nu = h\nu_0 + (\Delta m)\mu_{\text{B}}B$ avec $\Delta m \in \{0,
\pm1\}$~: trois positions seulement. (c) $\mu_{\text{B}}/h = \qty{14}{GHz/T}$, donc à $B = \qty{2}{T}$
l'écart vaut \qty{28}{GHz}, c'est-à-dire $\Delta\lambda =
\lambda^2\Delta\nu/c \approx \qty{23}{pm}$ --- résoluble depuis les
réseaux de Zeeman en 1896.
(d) Le spin propre de l'électron, avec son facteur anormal $g \approx
2$~: les motifs «~anormaux~» ont été les empreintes du spin avant
qu'il n'eût de nom (\cref{ch:b3:spin-two-level}).
\end{solution}

\begin{solution}{exo:b3:quantum-angular-momentum:8}
Dans les unités indiquées, $U_{\text{eff}} = -2/x + l(l+1)/x^2$. (a)
Pour $l = 1$~: minimum en $x = 2$ ($r = 2a_0$), de profondeur
$-E_{\text{I}}/2$. (b) $U_{\text{eff}} > 0$ pour $x < 1$~: en deçà
d'un rayon de Bohr, le mur l'emporte. (c) $\psi \sim r^l$ près de
l'origine~: seuls les états $l = 0$ ont une densité non nulle en
$r = 0$. (d) Capturer un électron exige une densité électronique
\emph{sur} le noyau~: les électrons $s$, seuls à n'être pas barrés par
le mur centrifuge, sont ceux qui sont avalés.
\end{solution}

\begin{solution}{exo:b3:quantum-angular-momentum:9}
(a) $\hat L_x = (\hat L_+ + \hat L_-)/2$ change $m$~: les éléments
diagonaux s'annulent. (b) Par symétrie les deux carrés transverses
sont égaux, et leur somme vaut $\langle\hat L^2 - \hat L_z^2\rangle$.
(c) À $m = l$~: $\Delta L_x\Delta L_y = l\hbar^2/2 = \tfrac\hbar2
\,l\hbar$~: c'est l'égalité --- l'état étiré est aussi aligné que
l'algèbre le permet. (d) Le cône~: $L_z$ bien défini, moyennes
transverses nulles, écarts transverses égaux --- un vecteur de
longueur et de projection définies, démocratiquement étalé en azimut.
\end{solution}

\begin{solution}{exo:b3:quantum-angular-momentum:10}
(a) Bilan d'énergie avec $E_J = BJ(J+1)$~: $J \to J + 1$ ajoute
$2B(J+1)$~; $J \to J - 1$ retranche $2BJ$. (b) $\Delta J = 0$ est
interdit, et les deux branches commencent à $\pm 2B$~: rien ne tombe
à la fréquence vibrationnelle pure. (c) Deux peignes d'espacement $2B$
encadrant un trou, d'intensités croissant jusqu'à $J_{\max}$ puis
décroissant. (d) La largeur de la bande est l'étendue peuplée des
peignes, qui croît comme $\sqrt{T}$~: ces ailes thermiquement peuplées
sont précisément là où une bande de serre saturée continue d'absorber.
\end{solution}

\begin{solution}{exo:b3:quantum-angular-momentum:11}
(a) $\|\hat J_\pm\ket{j,m}\|^2 = \bra{j,m}\hat J_\mp\hat J_\pm
\ket{j,m} = \hbar^2[j(j+1) - m(m\pm1)]$. (b) Grimper au-delà de $m = j$
donnerait $j(j+1) - j(j+1) = 0$ puis des normes négatives un pas plus
loin~: la chaîne doit contenir l'état supérieur annihilé. (c) Le haut
et le bas diffèrent d'un nombre entier de pas unité~: $2j \in
\mathbb N$~: $j = 0, \tfrac12, 1, \tfrac32, \dots$ (d) L'algèbre
n'utilise jamais de fonctions d'onde~; seule l'exigence que
$\eu^{\iu m\varphi}$ soit uniforme sur le cercle force $m$ entier ---
une condition qui lie le mouvement orbital et laisse les moments
cinétiques intrinsèques (de spin) libres d'être demi-entiers.
\end{solution}

\begin{solution}{exo:b3:quantum-angular-momentum:12}
(a) $n\hbar$ par unité de volume. (b) $\omega = n\hbar V/(\tfrac12 Ma^2)
= 2n\hbar/\rho a^2 = 2 \times \num{8.5e28} \times \num{1.055e-34}/
(7900 \times \num{e-4}) \approx \qty{2.3e-5}{rad/s}$ --- invisible en
un seul coup. (c) Renverser au rythme de la résonance du fil accumule
$\sim Q$ coups~: $\omega \sim \qty{2e-3}{rad/s}$, des oscillations
milliradian sur une période de quelques secondes --- visibles avec un
miroir et un faisceau lumineux, comme l'ont vu Einstein et de Haas.
(d) Le rapport gyromagnétique mesuré valait $e/m_{\text{e}}$, le
double de la valeur orbitale $e/2m_{\text{e}}$~: le magnétisme du fer
ne tient pas à des courants orbitaux mais au \emph{spin} de
l'électron, dont les chapitres suivants expliquent le $g \approx 2$.
\end{solution}

\begin{solution}{pb:b3:quantum-angular-momentum:1}
\textbf{1.} $E_J = BJ(J+1)$, dégénérescence $2J + 1$.
\textbf{2.} $I = \mu r_0^2 = \qty{1.46e-46}{kg.m^2}$~: $B =
\hbar^2/2I = \qty{3.8e-23}{J}$, soit $B/h = \qty{57.7}{GHz}$ ---
la constante mesurée.
\textbf{3.} $115$, $231$, \qty{346}{GHz}.
\textbf{4.} $E_{J+1} - E_J = 2B(J+1)$ croît linéairement~: deux raies
consécutives diffèrent de la constante $2B$. L'espacement donne $I$,
donc la longueur de liaison --- de la chimie structurale par radio.
\textbf{5.} $\lambda = c/\nu = \qty{2.6}{mm}$. La vapeur d'eau
atmosphérique absorbe avidement les ondes millimétriques~; d'où
l'Atacama, le Mauna Kea, le pôle Sud --- haut, froid, sec.
\textbf{6.} Le photon est une particule de spin $1$ portant $\hbar$~:
la conservation du moment cinétique total oblige le $J$ de la molécule
à changer d'une unité.
\textbf{7.} Sa distribution de charge est symétrique quel que soit
l'angle de rotation~: pas de dipôle oscillant, pas de raie dipolaire
--- un camouflage parfait.
\textbf{8.} $B_{\text{H}_2} = \hbar^2/2\mu r_0^2 \approx
\qty{7.6}{meV}$~; les faibles transitions quadrupolaires de cette
molécule symétrique exigent $\Delta J = 2$, au coût de $6B \approx
\qty{46}{meV}$ --- cinquante fois $k_{\text{B}}T$ à \qty{10}{K}~:
l'hydrogène froid ne peut absolument pas rayonner.
\textbf{9.} $2B = \qty{0.48}{meV} < k_{\text{B}}T = \qty{0.86}{meV}$~:
les collisions maintiennent les premiers barreaux peuplés ---
l'échelle fonctionne exactement là où vivent les nuages.
\textbf{10.} $n_1/n_0 = 3\,\eu^{-0.478/0.862} = 1.7$~: le niveau
émetteur est bien garni.
\textbf{11.} $J_{\max} \approx \sqrt{0.86/0.48} - 0.5 \approx 0.8$~:
la population culmine en $J = 1$.
\textbf{12.} CO livre les photons~; les convertir en masse totale de
gaz repose sur une abondance CO/H$_2$ supposée --- messager lumineux,
rançon étalonnée.
\textbf{13.} $\Delta\nu = \qty{115}{MHz}$~: $v = c\,\Delta\nu/\nu
\approx \qty{300}{km/s}$, en fuite (fréquence abaissée) --- des
vitesses orbitales galactiques.
\textbf{14.} $\pm\qty{2}{km/s}$ couvre \qty{1.5}{MHz}~; le mouvement
thermique à \qty{10}{K} ne donne que $\sim\qty{40}{kHz}$~: la largeur
est de la turbulence, non de la température --- les largeurs de raie
cartographient la météo interne d'un nuage.
\textbf{15.} Autour d'une masse centrale, $v \propto 1/\sqrt r$
devrait décroître (Kepler)~; la platitude mesurée signifie que la
masse enfermée continue de croître avec le rayon --- de la matière
invisible enveloppant le disque lumineux~: la matière noire.
\textbf{16.} Le rapport des populations, c'est-à-dire la température
d'excitation du gaz.
\textbf{17.} À \qty{57.6}{GHz}~: l'expansion cosmique a doublé toute
longueur d'onde en vol --- le décalage vers le rouge cosmologique sur
lequel revient le dernier chapitre de ce livre.
\textbf{18.} Flux de photons $\to$ luminosité CO (il faut la
distance) $\to$ nombre de CO (modèle d'excitation) $\to$ masse de
H$_2$ (le «~facteur X~» CO/H$_2$~: le maillon le plus faible, étalonné
localement et exporté avec prudence).
\textbf{19.} Un cœur à \qty{10}{K} n'émet rien qu'un télescope
optique puisse voir~; ses raies de rotation sont à la fois son seul
thermomètre, son seul tachymètre et sa seule balance.
\textbf{20.} Un thermomètre est sensible là où l'espacement vaut $\sim
k_{\text{B}}T$~: un espacement de \qty{5.5}{K}
répond fortement entre $5$ et $50$ K, tandis qu'un niveau optique à
\qty{2}{eV} ne serait peuplé par rien du tout.
\textbf{21.} La raie sature (milieu optiquement épais) et, pire, CO
gèle sur les grains de poussière dans le gaz le plus dense et le plus
froid~: les astronomes passent aux isotopologues plus rares
($^{13}$CO, C$^{18}$O) et à d'autres molécules.
\textbf{22.} Chaque espèce a besoin de sa propre densité et de sa
propre température pour briller~: ensemble, elles tomographient le
nuage --- enveloppes externes en CO, cœurs denses en NH$_3$ et HCN.
\textbf{23.} Une raie se voit par contraste sur le fond cosmique~; à
mesure que la température du gaz s'approche de \qty{2.7}{K}, son
excitation s'équilibre avec le fond et le contraste s'évanouit~:
aucun nuage plus froid ne se lit de cette façon.
\textbf{24.} $\lambda/20 = \qty{130}{\micro m}$ contre
$\qty{25}{nm}$ pour l'optique~: cinq mille fois plus tolérant ---
et c'est pourquoi des antennes millimétriques de douze mètres se
dressent en plein désert quand les miroirs optiques vivent sous
coupole.
\textbf{25.} Une échelle $BJ(J+1)$ avec $B/h = \qty{57.6}{GHz}$,
vivante à dix kelvins et lue à \qty{2.6}{mm}, mesure la masse, la
température, la turbulence et la rotation de l'univers froid --- et
ses courbes de rotation plates sont une pièce à conviction permanente
en faveur de la matière noire.
\end{solution}
